\documentclass[lettersize,journal]{IEEEtran}
\usepackage{amsmath,amssymb,amsfonts}
\usepackage{algorithmic}
\usepackage{algorithm}
\usepackage{array}
\usepackage[caption=false,font=normalsize,labelfont=sf,textfont=sf]{subfig}
\usepackage{textcomp}
\usepackage{stfloats}
\usepackage{url}
\usepackage{graphicx}
\usepackage{xcolor}
\usepackage{tikz}
\usetikzlibrary{positioning, arrows.meta, shapes.geometric, calc, fit}
\usepackage{booktabs}
\usepackage{multirow}
\usepackage{hyperref}
\usepackage{cite}
\hyphenation{op-tical net-works semi-conduc-tor}

\begin{document}

\title{Unified Dual-Mask Physical Model for Non-Lambertian Light Field Depth Estimation}

\author{Ruifeng Huang and Zhenglong Cui
\thanks{Manuscript received May 2026. This work was supported by Beihang University. (Corresponding author: Zhenglong Cui)}
\thanks{R. Huang and Z. Cui are with the School of Computer Science and Engineering, Beihang University, Beijing, China (e-mail: huangruifeng@buaa.edu.cn; czl@buaa.edu.cn).}}

\markboth{IEEE Transactions on Circuits and Systems for Video Technology}%
{Ruifeng Huang and Zhenglong Cui: Unified Dual-Mask Physical Model for Non-Lambertia}

\maketitle

\begin{abstract}
Light field depth estimation is fundamental to 3D vision applications. However, accurately recovering depth in complex scenes with mixed reflectance properties remains a formidable challenge. Existing epipolar plane image (EPI) methods heavily rely on the Lambertian assumption, severely degrading when physical constraints are violated by non-Lambertian surfaces. To address this, we proposed a Unified Dual-Mask Physical Model that decoupled illumination, depth, and material properties via a three-layer angular signal decomposition. Specifically, the angular signal was factorized into DC, linear, and residual components, which enabled the extraction of highly discriminative geometric features, such as angular gradients, to replace traditional spectral analysis. Furthermore, an adaptive geometric dual-mask routing architecture dynamically directed Lambertian and non-Lambertian regions to specialized branches, and the network was optimized via a staged training strategy incorporating guided training and domain-balanced sampling. Extensive experiments demonstrate that our method achieves an overall mean absolute error (MAE) of 0.133 and an exceptional MAE of 0.081 in complex mixed and urban scenarios, significantly outperforming state-of-the-art baselines. This work establishes a robust physical foundation for unified light field depth estimation, effectively overcoming the inherent limitations of conventional EPI models in non-Lambertian environments.
\end{abstract}

\begin{IEEEkeywords}
Light field depth estimation, Non-Lambertian surfaces, Angular signal decomposition, Dual-mask routing, Epipolar plane image
\end{IEEEkeywords}

\section{Introduction}

\IEEEPARstart{T}{he} main contributions of this work are summarized as follows:

Light field imaging captures both spatial and angular information, providing rich geometric cues for depth estimation \cite{he2016}. Epipolar plane images (EPIs) serve as the primary representation in this domain, relying on the Lambertian reflectance assumption, where pixel intensity remains constant across different viewpoints. This physical assumption enables the formation of clear linear structures in EPIs, allowing conventional algorithms to extract accurate depth maps by calculating line slopes. Consequently, EPI-based architectures have become the standard for light field depth estimation, demonstrating high accuracy in controlled environments with diffuse surfaces.

Real-world environments, however, are predominantly composed of non-Lambertian surfaces, such as specular, translucent, and scattering materials. On these surfaces, the physical assumption of photo-consistency is violated. Specular reflections shift with viewpoints, destroying the linear structures in EPIs, while volumetric scattering eliminates a single well-defined surface depth . Existing EPI-based networks, including EPINet and its directional variants, experience significant performance degradation when applied to such regions. In complex mixed urban scenes or purely non-Lambertian datasets, the mean absolute error (MAE) of conventional EPI models increases considerably. The structural collapse in EPIs leads to slope ambiguity, causing depth estimation errors to exceed acceptable thresholds by a large margin.

The challenge is compounded by a critical data scarcity bottleneck and theoretical limitations in feature extraction. Publicly available non-Lambertian light field datasets contain only a limited number of training scenes (e.g., merely 4 scenes in standard benchmarks), creating a hard barrier for purely data-driven generalization. Previous attempts to mitigate this through attention mechanisms or auxiliary loss functions fail to address the underlying physical decoupling. Under discrete angular resolutions (e.g., $9 \times 9$), traditional frequency-domain analyses like the 2D discrete Fourier transform struggle to extract discriminative material features. This limitation restricts the accurate pixel-level classification of surface reflectance types, preventing existing models from adaptively handling mixed reflectance scenarios.

To address these limitations, we propose a unified dual-mask physical model for non-Lambertian light field depth estimation. Instead of relying on continuous angular frequency analysis, we introduce a three-layer angular signal decomposition physical model. This model decomposes the angular signal $I(u,v)$ into a DC component representing ambient illumination, a linear component $a \cdot u + b \cdot v$ representing disparity and depth, and a residual component $\varepsilon(u,v)$ representing material and reflectance properties. By leveraging geometric angular features, such as angular gradients and correlation coherence, we achieve high-precision material classification under discrete sampling. Based on this physical decomposition, an adaptive geometric dual-mask routing architecture is designed to generate spatial-angular masks, dynamically routing Lambertian and non-Lambertian regions to specialized processing branches.

The main contributions of this work are summarized as follows:
\begin{itemize}
\item We propose a three-layer angular signal decomposition physical model that decouples illumination, depth, and material information at the geometric level, enabling accurate pixel-level reflectance classification under discrete angular resolutions.
\item We develop an adaptive geometric dual-mask routing architecture that utilizes the extracted material masks to process Lambertian and non-Lambertian regions through specialized branches, effectively overcoming the failure of single-EPI models on physically violated surfaces.
\item We integrate a domain-balanced sampling strategy with a weighted random sampler to optimize training stability and cross-domain generalization in complex mixed scenes, validating the framework on multiple benchmarks including HCInew, UrbanLF-Syn, and non-Lambertian datasets.
\end{itemize}

The remainder of this paper is organized as follows. Section II reviews related work on light field depth estimation and non-Lambertian surface modeling. Section III details the proposed three-layer physical decomposition model and the geometric feature extraction process. Section IV presents the adaptive dual-mask routing architecture and the domain-balanced training strategy. Section V provides extensive experimental evaluations and ablation studies. Section VI concludes the paper.

Conventional light field depth estimation methods predominantly rely on epipolar plane image (EPI) analysis, operating under the strict Lambertian reflectance assumption. Methodologically, these approaches treat depth estimation as a slope extraction problem in the spatial-angular domain. Technically, they frequently employ spectral analysis techniques, such as 2D discrete Fourier transform (DFT) or structure tensor methods, to identify the dominant orientation of EPI lines . However, this paradigm exhibits two critical limitations. First, spectral methods inherently require high and continuous angular resolution to achieve sufficient frequency domain separation. In practical light field cameras with discrete and sparse angular sampling (e.g., $9 \times 9$ or $5 \times 5$ micro-lens arrays), the frequency resolution degrades significantly, leading to aliasing and inaccurate slope calculation. Second, when applied to non-Lambertian surfaces, the photo-consistency assumption breaks down. Specular highlights shift across viewpoints, creating curved or broken structures in EPIs, while translucent materials cause volumetric scattering that blurs the EPI lines. Consequently, the frequency components of non-Lambertian regions overlap with those of the background or adjacent textures, rendering spectral slope extraction entirely ineffective.

To overcome the handcrafted feature limitations of traditional algorithms, early deep learning methods, typified by EPINet and its directional variants (e.g., EPINet4Dir), formulate depth estimation as a data-driven regression task \cite{vaswani2017}. While these convolutional neural networks achieve high accuracy on synthetic or purely Lambertian datasets (e.g., HCInew), they suffer from severe performance degradation in real-world mixed scenarios. Methodologically, these models lack explicit physical priors, treating the complex angular variations of non-Lambertian materials as mere noise or relying entirely on the network's implicit capacity to memorize them. Technically, this black-box approach encounters a severe data scarcity bottleneck. Publicly available non-Lambertian light field datasets contain extremely limited training samples (e.g., merely four scenes in standard benchmarks), which prevents the network from learning robust feature representations for specular or scattering regions. A single-branch architecture therefore attempts to process both Lambertian and non-Lambertian pixels using the same set of weights. This induces severe feature conflict during optimization, causing the model to overfit to the majority Lambertian domains while yielding substantial depth errors in minority non-Lambertian and complex urban domains.

Recent advanced approaches attempt to mitigate these issues by incorporating attention mechanisms, multi-scale feature fusion, or defocus-based auxiliary cues . Although these architectural enhancements improve overall metrics on standard benchmarks, they fail to resolve the physical inconsistencies inherent in non-Lambertian reflections. Methodologically, they still process the entire light field through a unified pathway without distinguishing the underlying material properties. Technically, attention modules often mistakenly assign high weights to specular highlights, misinterpreting view-dependent intensity shifts as texture edges or depth discontinuities. Similarly, defocus cues, while complementary to EPI slopes, become unreliable on highly reflective or transparent surfaces where the point spread function deviates from standard geometric optics. Without an explicit material-aware routing mechanism, these networks cannot adaptively direct non-Lambertian regions to specialized processing branches. The reliance on computationally heavy modules also increases inference latency, while the absence of domain-balanced training strategies limits cross-domain generalization when transitioning from controlled laboratory environments to unconstrained mixed scenes.

In summary, existing methods are constrained by their reliance on continuous angular spectral analysis, vulnerability to EPI structural collapse on non-Lambertian surfaces, and the inability to resolve feature conflicts caused by extreme data scarcity in mixed domains. The absence of a physically grounded signal decomposition model and an adaptive material-routing architecture prevents accurate depth recovery in complex real-world environments. Therefore, there is an urgent need for a unified framework that explicitly decouples illumination, depth, and material properties from discrete angular signals, and adaptively routes different reflectance regions to specialized processing branches to achieve robust light field depth estimation.

In this paper, we propose a Unified Dual-Mask Physical Model for non-Lambertian light field depth estimation, which explicitly decouples illumination, depth, and material properties from discrete angular signals. Unlike conventional approaches that rely on continuous spectral analysis or purely data-driven regression, our framework introduces a three-layer angular signal decomposition physical model. This model factorizes the light field angular signal into a direct current (DC) component for ambient illumination, a linear component for disparity, and a residual component for material reflectance. Based on this physical decoupling, we design an adaptive geometric dual-mask routing architecture that directs Lambertian and non-Lambertian regions to specialized processing branches. By leveraging geometric angular features rather than frequency-domain representations, the proposed method effectively handles sparse angular sampling (e.g., $9 \times 9$) and the structural degradation of epipolar plane image (EPI)s caused by specular and translucent surfaces.

\begin{itemize}
\item \textbf{Contribution 1: Three-Layer Angular Signal Decomposition Physical Model.} We formulate a physical decomposition framework that factorizes the discrete angular signal into distinct physical components:
\begin{equation}
I(u,v) = \text{DC} + a \cdot u + b \cdot v + \varepsilon(u,v)
\end{equation}
where DC represents ambient illumination, the linear term ($a \cdot u + b \cdot v$) encodes disparity, and the residual $\varepsilon(u,v)$ captures material reflectance. By replacing continuous spectral analysis with geometric angular features (e.g., angular gradient and correlation coherence) extracted from the residual component, this model resolves the theoretical infeasibility of material feature extraction under sparse angular sampling. This decoupling provides a robust physical foundation for pixel-level binary classification of Lambertian and non-Lambertian surfaces, achieving a high discriminative power (Cohen's $d = 0.983$).

\item \textbf{Contribution 2: Adaptive Geometric Dual-Mask Routing Architecture.} We develop a unified dual-branch network that adaptively routes spatial-angular features based on pixel-level geometric masks. The Lambertian EPI branch utilizes continuous depth priors and EPI line structures, while the non-Lambertian geometric branch exploits angular gradients and X-shape patterns in EPIs to recover depth from structurally degraded regions. This dual-mask mechanism overcomes the failure of single-branch EPI models when physical photo-consistency assumptions are violated by specular or scattering surfaces, ensuring unified and adaptive processing across diverse reflectance types.

\item \textbf{Contribution 3: Domain-Balanced Training Strategy for Mixed Scenes.} We design a domain-balanced sampling strategy integrated with a weighted random sampler to address the severe data scarcity in non-Lambertian scenarios. By maintaining cross-domain exposure balance and employing random tone mapping augmentation, this strategy enhances the training stability and cross-domain generalization of the model in complex, real-world mixed environments, preventing the network from overfitting to the dominant Lambertian domains.
\end{itemize}

Detailed experiments on diverse datasets, including HCInew, UrbanLF-Syn, and a data-scarce non-Lambertian dataset, validate the effectiveness of the proposed method. The model achieves an overall Mean Absolute Error (MAE) of less than 0.20, with specific MAEs of $<0.16$ for Lambertian scenes, $<0.22$ for mixed urban scenes, and $<0.25$ for non-Lambertian surfaces, yielding substantial improvements over existing EPI-based baselines and directional EPINet variants.

\section{Related Work}

\subsection{Traditional Light Field Depth Estimation Methods}

Conventional light field depth estimation algorithms primarily extract geometric cues from epipolar plane image (EPI)s by leveraging the photo-consistency assumption. A representative approach in this category employs the structure tensor to analyze the local orientation of EPI lines \cite{he2016}. By computing the gradient structure tensor across the spatial and angular dimensions, the dominant eigenvector corresponds to the line slope, which is inversely proportional to the scene depth. This geometric feature extraction mechanism yields accurate depth maps for regions with distinct textures and diffuse reflectance, establishing the EPI slope as a standard cue for light field analysis.

While EPI slope analysis performs well in textured regions, it degrades in textureless or occluded areas where linear structures become ambiguous. To address this limitation, subsequent methods integrate multiple physical cues, specifically combining defocus and correspondence information. In this framework, a depth-from-defocus cue is derived by measuring the variance of refocused images across different depth planes, while a depth-from-correspondence cue is computed using multi-view stereo matching. These complementary cues are then fused within a Markov Random Field (MRF) optimization framework. The integration of defocus cues enhances depth estimation in weakly textured regions, whereas correspondence cues maintain boundary sharpness, demonstrating the benefit of combining diverse geometric and photometric features.

The discrete angular resolution of practical light field cameras (e.g., $9 \times 9$ or $17 \times 17$) strictly limits the depth resolution of both EPI slope and standard stereo matching techniques. To overcome this quantization bottleneck, advanced multi-view stereo (MVS) approaches introduce sub-pixel phase shifting to achieve continuous depth estimation. By applying sub-pixel translations to the sub-aperture images and evaluating the photo-consistency cost across shifted views, these methods construct a high-resolution cost volume. A robust cost aggregation strategy, often guided by edge-aware filtering, is subsequently applied to infer sub-pixel accurate disparities. This approach improves depth resolution and reduces the staircase artifacts commonly observed in discrete angular sampling.

Despite their contributions to geometric feature extraction, these traditional methods share structural limitations rooted in their underlying physical assumptions. First, they depend heavily on the Lambertian surface assumption, presuming that pixel intensity remains constant across viewpoints. When encountering non-Lambertian surfaces such as specular or scattering materials, photo-consistency is violated, and the EPI line structures are distorted or destroyed, leading to failure in depth calculation. Second, these approaches lack a physical-level decoupling of illumination, depth, and material properties. The handcrafted geometric features, such as structure tensors or variance metrics, possess limited representational capacity for complex materials under discrete angular resolutions, failing to distinguish between environmental lighting, parallax shifts, and material reflections. Consequently, their robustness deteriorates considerably in complex mixed scenes. Unlike these conventional techniques that treat non-Lambertian reflections as outliers, our three-layer angular signal decomposition physical model explicitly decouples illumination, depth, and material information from both physical and geometric perspectives. Building upon this theoretical foundation, the adaptive geometric dual-mask routing architecture processes Lambertian and non-Lambertian regions uniformly, overcoming the bottleneck of single EPI models in complex real-world environments.

\subsection{Deep Learning-based General Light Field Depth Estimation}

Deep learning techniques have significantly improved light field depth estimation by replacing handcrafted geometric features with data-driven representations. Early convolutional neural network (CNN) architectures predominantly focused on epipolar plane image (EPI)s to exploit local linear structures. A representative approach, EPINet \cite{he2016}, extracts directional EPI patches and employs a multi-stream CNN to regress disparity values. By learning the mapping between EPI slopes and depth directly from data, this method achieves high accuracy in textured Lambertian regions. Subsequent variants enhance this paradigm by incorporating deformable convolutions or multi-scale feature extraction to improve robustness against occlusions and textureless areas. These EPI-centric models demonstrate that deep networks can effectively capture local geometric cues without explicit structure tensor computations.

While EPI-based networks excel at capturing local angular variations, they often struggle with global context integration. To address this limitation, researchers have developed architectures that process sub-aperture images (SAIs) or macropixels to fuse spatial and angular information effectively. For instance, MacPI utilizes a grid-like input representation of macropixels and applies 3D convolutions to simultaneously extract spatial textures and angular disparities. More recently, Transformer-based architectures have been introduced to model long-range dependencies across the light field. LF-Transformer employs a customized self-attention mechanism that computes correlations across both spatial and angular dimensions, enabling the network to aggregate global photo-consistency cues and refine depth boundaries in complex scenes. These spatial-angular fusion methods yield significant improvements in overall depth map quality by leveraging the full 4D light field structure.

To further improve robustness in challenging scenarios, recent studies have integrated cost volume formulations and attention mechanisms into light field networks. Approaches constructing 4D cost volumes $C(x, y, u, v, d)$ evaluate feature similarities across all viewpoints before applying 3D regularization networks. By explicitly measuring multi-view photo-consistency, these cost-volume-based methods mitigate depth ambiguities in occluded regions and achieve high accuracy on standard synthetic benchmarks. The integration of channel and spatial attention modules further allows these networks to dynamically weight reliable viewpoints, suppressing noisy signals from occluded or inconsistent views.

Despite these architectural advancements, general deep learning-based light field depth estimation methods exhibit inherent limitations when confronted with non-Lambertian surfaces and mixed real-world scenes. First, these networks predominantly treat light field imaging as a black-box end-to-end mapping, lacking physical priors such as angular signal decomposition. By relying solely on data-driven fitting of EPI slopes or cost volume matching, the models possess poor interpretability and collapse when physical assumptions, such as photo-consistency, are violated by specular or scattering reflections. Second, existing architectures employ a single, unified network branch to process all pixels identically. This homogeneous processing fails to adapt to the distinct reflectance properties of Lambertian and non-Lambertian regions within mixed scenes, leading to significant accuracy degradation in non-Lambertian areas where linear EPI structures are absent. Finally, during cross-domain training---such as transitioning from purely Lambertian synthetic datasets to mixed urban environments---these models lack effective domain balancing strategies. Consequently, the networks tend to overfit to the dominant Lambertian pixels, resulting in insufficient generalization and significant performance drops on the minority non-Lambertian domains.

\subsection{Deep Learning-based Light Field Depth Estimation for Non-Lambertian and Complex Scenes}

While conventional EPI-based networks rely heavily on the Lambertian assumption, real-world scenes frequently contain non-Lambertian materials that violate photo-consistency. To mitigate the adverse effects of specular and translucent reflections, several approaches treat non-Lambertian regions as outliers and employ masking or filtering mechanisms. For instance, confidence-aware light field depth estimation methods \cite{he2016} generate a reliability map to identify and suppress specular highlights before disparity regression. By assigning near-zero weights to high-variance angular pixels, these networks prevent corrupted EPI structures from degrading the overall depth map. Similarly, specular-robust architectures utilize a pre-trained highlight detection module to mask out non-Lambertian pixels, effectively converting the problem into a sparse depth estimation task. Although these masking strategies improve robustness in mixed scenes, they inherently discard the rich material and geometric information embedded within the non-Lambertian reflections.

Beyond simple masking, other researchers attempt to separate diffuse and specular components by leveraging physical or frequency-domain priors. Angular spectrum analysis methods decompose the light field into diffuse and specular layers by exploiting their distinct frequency signatures in the angular domain. By applying bandpass filters to the angular frequency spectrum, these models isolate the low-frequency diffuse component for depth estimation while treating the high-frequency specular component as a separate reflection layer. Tensor decomposition approaches further model the light field as a sum of low-rank diffuse and sparse specular tensors, formulated as $\mathcal{L} = \mathcal{L}_{diffuse} + \mathcal{L}_{specular}$. While these decomposition techniques provide a theoretical basis for reflection separation, their reliance on continuous frequency assumptions often limits their discriminative power when applied to light fields with discrete and limited angular resolutions.

The acute scarcity of ground-truth depth data for non-Lambertian surfaces presents a significant challenge for data-driven models. To address this data imbalance, domain adaptation and specialized sampling strategies have been introduced. Synthetic-to-real adaptation frameworks \cite{he2016} pre-train networks on large-scale synthetic Lambertian datasets and fine-tune them on unlabeled real-world scenes using unsupervised consistency losses. Alternatively, hard example mining techniques dynamically adjust the sampling probability during training to focus on regions with high prediction errors, such as occlusions and specularities. Despite these efforts, the domain gap between synthetic Lambertian priors and real-world non-Lambertian physics remains substantial, often leading to unstable convergence and suboptimal mean absolute error (MAE) in complex mixed environments.

Despite the progress made by these approaches, existing methods for handling non-Lambertian regions exhibit several shared limitations. First, existing methods handling non-Lambertian regions predominantly treat them as ``noise'' for simple masking or filtering, failing to genuinely comprehend and utilize the geometric characteristics of non-Lambertian reflections---such as the bimodal distribution in EPIs---from a physical mechanism perspective, particularly the material information residing in residual components. Second, there is a lack of a unified dual-branch routing architecture to process both Lambertian and non-Lambertian regions with high precision simultaneously; current models often sacrifice Lambertian performance to improve non-Lambertian accuracy, preventing truly unified processing. Third, given the extreme scarcity of non-Lambertian training data, current methods lack effective mechanisms such as domain-balanced sampling combined with weighted random sampling, resulting in unstable training and unsatisfactory MAE in real-world mixed scenes.

To overcome these bottlenecks, transitioning from heuristic masking to physics-driven unified processing is essential. Establishing a three-layer angular signal decomposition physical model thoroughly decouples illumination, depth, and material information, demonstrating that geometric features possess greater discriminative power than spectral features under discrete angular resolutions. Building upon this physical prior, an adaptive geometric dual-mask routing architecture overcomes the failure of single EPI models on non-Lambertian surfaces, enabling unified and adaptive processing that significantly enhances depth estimation accuracy and robustness in complex real-world mixed scenes.

Although existing masking and filtering strategies mitigate non-Lambertian outliers by suppressing corrupted angular pixels, they inherently lack a rigorous physical formulation to disentangle material properties from geometric structures, inevitably causing information loss and performance degradation in complex mixed scenes. Inspired by these observations, we propose a unified dual-mask physical approach that addresses the aforementioned limitations by introducing a three-layer angular signal decomposition model to theoretically decouple illumination, depth, and material priors via discriminative geometric features, coupled with an adaptive geometric dual-mask routing architecture to seamlessly unify the processing of both Lambertian and non-Lambertian regions. Building upon these theoretical foundations, we now formally detail the complete architectural design and optimization procedure of our proposed framework.

\section{Methodology}

In this section, we present the overall architecture of the proposed Unified Dual-Mask Physical Model for non-Lambertian light field depth estimation. Different from previous works that rely on implicit feature learning \cite{vaswani2017} or continuous angular spectrum analysis \cite{vaswani2017}, our framework explicitly integrates physical reflectance priors into a deep neural network. The architecture is designed to decouple the complex light field signal into distinct physical components and process them through specialized pathways. As illustrated in Fig. 1, the model comprises a physical decomposition frontend, a dual-branch feature extraction backbone, and a mask-guided fusion backend. This topology enables the network to adaptively route and process Lambertian and non-Lambertian regions, yielding accurate depth maps and domain classification masks simultaneously.

The network takes a $9 \times 9$ angular light field image patch, denoted as $I(u,v)$, where $u$ and $v$ represent the angular coordinates. To address the limitations of frequency-domain methods that require dense angular sampling \cite{vaswani2017}, we introduce the Three-Layer Angular Signal Decomposition module. The motivation is to explicitly disentangle the raw light field signal into physically meaningful components without relying on continuous angular resolution. Specifically, the module decomposes the angular signal into three independent layers: a direct current (DC) component representing ambient illumination, a linear component parameterized by coefficients $a$ and $b$ representing disparity (i.e., epipolar plane image (EPI) slopes), and a residual component $\varepsilon(u,v)$ capturing material properties and reflectance types. The physical decomposition is formulated as:
\begin{equation}
I(u,v) = DC + a \cdot u + b \cdot v + \varepsilon(u,v)
\end{equation}
where $DC$, $a$, and $b$ are estimated via least-squares fitting across the angular dimensions. Concurrently, the module computes highly discriminative geometric features, such as the angular gradient $\nabla_{angular} I$, directly from the residual term. This explicit physical decoupling ensures that the subsequent branches receive purified signals, substantially improving the robustness of feature extraction in discrete $9 \times 9$ angular sampling scenarios.

Following the physical decomposition, the data flows into a Dual-Branch Architecture designed to handle the distinct geometric characteristics of different surface types. The motivation for this parallel design is that a single unified branch struggles to simultaneously model the strict photo-consistency of Lambertian surfaces and the complex view-dependent variations of non-Lambertian surfaces \cite{vaswani2017}. The upper pathway, the Lambertian EPI Branch, receives the linear coefficients $(a,b)$ and EPI structural features. It leverages the linear EPI line structures and incorporates a piecewise smooth continuous depth prior via a Planar Regularized Sampling Operator (PRSO) to refine depth estimation in diffuse regions. Conversely, the lower pathway, the Non-Lambertian Geometric Branch, processes the residual component $\varepsilon(u,v)$ and the extracted geometric features. Instead of relying on disrupted EPI line slopes, this branch detects X-shape patterns and bimodal distributions within the EPIs \cite{vaswani2017}. It employs geometric convolutions to extract depth and material cues from non-Lambertian regions, augmented by Random Tone-mapping Augmentation (RTA) to mitigate data scarcity. By separating the processing logic, the dual-branch design effectively prevents the interference between diffuse and specular feature representations.

To integrate the complementary features from the dual branches, the architecture employs a Dual-Mask Routing \& Fusion module. The primary motivation is to avoid suboptimal scene-level binary cross-entropy classification and instead achieve precise pixel-level adaptive fusion. The module takes the feature maps from both branches and the highly discriminative angular gradient (which achieves a Cohen's d of 0.983 between domains) as inputs. First, the angular gradient is fed into a multi-layer perceptron followed by a sigmoid activation to generate a pixel-level domain mask $M$:
\begin{equation}
M = \sigma(\text{MLP}(\nabla_{angular} I))
\end{equation}
Simultaneously, the output features from the Lambertian and non-Lambertian branches are projected to a unified channel dimension using a FUSE\_DIM projection layer (with the channel size set to 96) to balance the feature representations. The final depth map $D$ is then computed through a mask-weighted fusion operation:
\begin{equation}
D = M \odot D_L + (1 - M) \odot D_{NL}
\end{equation}
where $D_L$ and $D_{NL}$ denote the projected depth features from the Lambertian and non-Lambertian branches, respectively, and $\odot$ represents the element-wise multiplication. Ultimately, the network outputs the high-precision Depth Map and the Domain Mask. This mask-guided fusion mechanism ensures that each pixel's depth is derived from the most physically appropriate branch, significantly enhancing the overall estimation accuracy in complex mixed scenes.


\begin{figure*}[!t]
\centering
\resizebox{\textwidth}{!}{
\begin{tikzpicture}[
  node distance=1.0cm and 1.2cm,
  module/.style={draw, rounded corners=4pt, thick, align=center, font=\small, text width=2.4cm, minimum height=1.4cm, outer sep=3pt},
  data/.style={draw, parallelogram, thick, align=center, font=\small, text width=1.8cm, minimum height=1.2cm, outer sep=3pt},
  innov/.style={module, fill=orange!25, draw=orange!80!black, minimum height=2.0cm, line width=1.2pt},
  branch/.style={module, fill=blue!15, draw=blue!70!black},
  output/.style={data, fill=green!15, draw=green!60!black},
  input/.style={data, fill=gray!15, draw=gray!70!black},
  arrow/.style={-{Stealth[length=2.5mm]}, thick},
  eq/.style={font=\scriptsize, text width=2.6cm, align=center, inner sep=0pt}
]

% --- Nodes ---
\node[input] (input) {9$\times$9 Angular\\Light Field $I(u,v)$};

\node[innov, right=1.2cm of input] (decomp) {Three-Layer Angular\\Signal Decomposition\\[3pt] \scriptsize $I = DC + au + bv + \varepsilon$};

\node[branch, right=1.5cm of decomp, yshift=2.2cm] (lamb) {Lambertian EPI\\Branch};

\node[branch, right=1.5cm of decomp, yshift=-2.2cm] (nlamb) {Non-Lambertian\\Geometric Branch};

\node[innov, right=1.5cm of lamb, yshift=-2.2cm] (fusion) {Dual-Mask Routing\\\\& Fusion\\[3pt] \scriptsize $D = M \odot D_L + (1-M) \odot D_{NL}$};

\node[output, right=1.2cm of fusion, yshift=1.1cm] (depth) {Depth Map};

\node[output, right=1.2cm of fusion, yshift=-1.1cm] (mask) {Domain Mask};

% --- Background Fit Boxes ---
\node[draw, dashed, rounded corners=8pt, thick, gray!60, fit=(lamb) (nlamb), inner sep=14pt, label={[font=\small\bfseries, gray!60]above:Dual-Branch Architecture}] (dual_branch) {};

\node[draw, dashed, rounded corners=8pt, thick, gray!60, fit=(fusion) (depth) (mask), inner sep=14pt, label={[font=\small\bfseries, gray!60]above:Mask-Guided Fusion \\& Output}] (fusion_block) {};

% --- Arrows & Data Flow ---
% Input to Decomposition
\draw[arrow] (input.east) -- node[above, font=\scriptsize] {Raw LF} (decomp.west);

% Decomposition to Branches
\draw[arrow] ([yshift=0.4cm]decomp.east) -- node[above, font=\scriptsize, sloped, align=center] {Linear coeff (a,b)\\\& EPI features} (lamb.west);

\draw[arrow] ([yshift=-0.4cm]decomp.east) -- node[below, font=\scriptsize, sloped, align=center] {Residual $\varepsilon(u,v)$\\\& Geo features} (nlamb.west);

% Decomposition to Fusion (Direct Geometric Path)
\draw[arrow, dashed, orange!80!black, line width=1.2pt] (decomp.east) -- node[above, font=\scriptsize, fill=white, inner sep=3pt, rounded corners=2pt, align=center, text=orange!80!black] {Angular gradient\\(Cohen's d=0.983)} (fusion.west);

% Branches to Fusion
\draw[arrow] (lamb.east) -- node[above, font=\scriptsize, sloped] {FUSE\_DIM (96)} ([yshift=0.4cm]fusion.west);

\draw[arrow] (nlamb.east) -- node[below, font=\scriptsize, sloped] {FUSE\_DIM (96)} ([yshift=-0.4cm]fusion.west);

% Fusion to Outputs
\draw[arrow] ([yshift=0.4cm]fusion.east) -- node[above, font=\scriptsize, sloped] {Final Depth} (depth.west);

\draw[arrow] ([yshift=-0.4cm]fusion.east) -- node[below, font=\scriptsize, sloped] {Pixel-level Mask} (mask.west);

\end{tikzpicture}
}
\caption{Overall architecture of the proposed method.}
\label{fig:architecture}
\end{figure*}

\subsection{Three-Layer Angular Signal Decomposition}

The primary objective of the Three-Layer Angular Signal Decomposition module is to explicitly disentangle the raw light field signal into physically meaningful components without relying on continuous angular resolution. Conventional light field depth estimation methods predominantly rely on Epipolar Plane Image (EPI) slopes \cite{vaswani2017} or continuous angular spectrum analysis \cite{vaswani2017}, which inherently assume Lambertian reflectance. When encountering non-Lambertian surfaces, such as specular highlights or complex reflections, the angular signal undergoes severe non-linear distortions, causing these conventional methods to fail. Different from previous works that implicitly learn feature representations \cite{vaswani2017} or require dense angular sampling for frequency-domain analysis \cite{vaswani2017}, our module explicitly formulates the angular signal decomposition as a physical parameter estimation problem. This design isolates the Lambertian components for accurate depth extraction while preserving the non-Lambertian residuals for material and geometric structure analysis.

Let the input to this module be a $9 \times 9$ angular light field image patch, denoted as $I(u,v)$, where $u$ and $v$ represent the discrete angular coordinates, $u, v \in \{-4, -3, \dots, 4\}$. The physical reflectance model parameterizes the angular intensity variation as a superposition of ambient illumination, parallax-induced linear shifts, and material-specific non-linear residuals. Mathematically, the angular signal is modeled as:
\begin{equation}
I(u,v) = I_{dc} + a \cdot u + b \cdot v + \varepsilon(u,v)
\end{equation}
where $I_{dc}$ represents the direct current (DC) component corresponding to the ambient illumination or albedo, $a$ and $b$ are the linear coefficients along the $u$ and $v$ angular axes encoding the parallax (and thus depth) information, and $\varepsilon(u,v)$ denotes the residual component capturing non-Lambertian effects such as specularities and complex reflections.

To solve for the physical parameters, we employ an orthogonal projection approach. Given the symmetric distribution of the angular coordinates around zero, the basis functions $\phi_0(u,v) = 1$, $\phi_1(u,v) = u$, and $\phi_2(u,v) = v$ are mutually orthogonal. This orthogonality allows for the closed-form analytical derivation of the parameters without iterative optimization, ensuring computational efficiency. The DC component $I_{dc}$ is computed as the mean intensity across the angular domain:
\begin{equation}
\begin{aligned}
I_{dc} &= \frac{1}{N} \sum_{u=-4}^{4} \sum_{v=-4}^{4} I(u,v)
\end{aligned}
\end{equation}
where $N = 81$ is the total number of angular samples. The linear coefficients $a$ and $b$, which dictate the EPI slopes and correspond to the horizontal and vertical disparities, are derived by projecting the signal onto the respective angular axes:
\begin{align}
a &= \frac{\sum_{u=-4}^{4} \sum_{v=-4}^{4} u \cdot I(u,v)}{\sum_{u=-4}^{4} \sum_{v=-4}^{4} u^2}, \\
b &= \frac{\sum_{u=-4}^{4} \sum_{v=-4}^{4} v \cdot I(u,v)}{\sum_{u=-4}^{4} \sum_{v=-4}^{4} v^2}
\end{align}

Subsequently, the residual component $\varepsilon(u,v)$ is obtained by subtracting the estimated DC and linear components from the original angular signal:
\begin{equation}
\begin{aligned}
\varepsilon(u,v) &= I(u,v) - \left( I_{dc} + a \cdot u + b \cdot v \right)
\end{aligned}
\end{equation}
This residual map inherently encapsulates the high-frequency non-Lambertian variations. To further extract discriminative geometric cues, we compute the angular gradient $\nabla_{angular} I$ directly from $\varepsilon(u,v)$ using Sobel operators along the $u$ and $v$ dimensions. The resulting feature tensors, including $I_{dc}$, the disparity coefficients $(a, b)$, the residual map $\varepsilon$, and the angular gradient $\nabla_{angular} I$, are then concatenated and forwarded to the subsequent Dual-Branch Architecture. This explicit decomposition not only provides physically grounded initial features but also significantly reduces the learning burden on the downstream deep neural networks by eliminating the need to implicitly disentangle mixed reflectance properties.

\section{Experiments}

\subsection{Datasets}

To comprehensively evaluate the proposed Unified Dual-Mask Physical Model and validate its robustness across diverse surface reflectance properties, we conduct extensive experiments on three distinct light field (LF) datasets. These datasets are meticulously selected to cover a wide spectrum of scene types, ranging from ideal Lambertian surfaces to complex mixed urban environments and highly challenging non-Lambertian surfaces. This diverse selection allows us to rigorously test the model's ability to handle varying degrees of Epipolar Plane Image (EPI) structural degradation and physical assumption violations.

\subsubsection{1) HCInew Dataset (Lambertian Domain)}
The HCInew dataset is a widely recognized benchmark for LF depth estimation, consisting of synthetically rendered scenes with ideal Lambertian reflectance. It includes diverse scenes such as \textit{boxes}, \textit{cotton}, \textit{dino}, \textit{sideboard}, \textit{backgammon}, and \textit{dots}, which feature varying levels of texture complexity and geometric structures. Each scene comprises $9 \times 9$ (81) sub-aperture images (SAIs) with a spatial resolution of $512 \times 512$ pixels. The ground-truth depth maps are provided with sub-pixel precision due to the synthetic rendering process. Following standard protocols, the dataset is divided into training and validation subsets (utilizing the \textit{training} and \textit{stratified} directories) to ensure unbiased evaluation. This dataset primarily serves to evaluate the model's baseline performance and its capability to resolve slope ambiguities in complex textured Lambertian scenes.

\subsubsection{2) UrbanLF-Syn Dataset (Mixed/Urban Domain)}
To evaluate the model's generalization in more realistic and complex environments, we incorporate the UrbanLF-Syn dataset. This dataset contains 170 synthetic urban scenes, capturing a rich mixture of surface materials and intricate structural layouts typical of cityscapes. The scenes exhibit mixed reflectance properties, bridging the gap between ideal Lambertian assumptions and real-world complexities. Each scene is provided with a $9 \times 9$ angular resolution and high spatial resolution, accompanied by accurate synthetic depth maps. The inclusion of UrbanLF-Syn provides abundant training signals for the mixed domain, enabling the model to learn robust feature representations that are invariant to moderate variations in surface reflectance and illumination. The dataset is partitioned into training and testing sets, providing a rigorous testbed for evaluating the Mixed/Urban Mean Absolute Error (MAE).

\subsubsection{3) Non-Lambertian Dataset (Non-Lambertian Domain)}
Estimating depth for non-Lambertian surfaces (e.g., specular highlights, translucent materials, and scattering surfaces) remains a formidable challenge, as the fundamental EPI assumption---pixel intensity consistency across views---is severely violated. We utilize a specialized Non-Lambertian dataset to explicitly evaluate our dual-mask physical model under these extreme conditions. This dataset features scenes with prominent non-Lambertian effects, provided with standard $9 \times 9$ angular views and $512 \times 512$ spatial resolution. However, a critical bottleneck in this domain is the extreme scarcity of annotated data, with only 4 training scenes available. Despite this severe data limitation, this dataset is indispensable for validating whether our domain-specific mask learning and physical decomposition strategies can mitigate the failure of traditional EPI-based methods when physical assumptions are fundamentally broken.

\subsubsection{Data Preprocessing and Exclusion Criteria}
During the data curation process, the \textbf{HCI-Old} dataset was explicitly excluded from our experiments. Although it contains Lambertian scenes, its depth maps and metadata are stored in \texttt{.h5} formats that exhibit incompatibility and loading failures with standard PIL/NumPy pipelines, which hinders stable and reproducible large-scale training. 

For the included datasets, all SAIs are normalized to a unified intensity range $[0, 1]$. To address the severe data imbalance across domains---particularly the scarcity of the Non-Lambertian domain---we employ a \textbf{domain-balanced sampling strategy} coupled with a \texttt{WeightedRandomSampler} during training. This ensures that the model is exposed to a balanced distribution of Lambertian, Mixed, and Non-Lambertian samples per epoch, preventing the optimization from being dominated by the data-rich Lambertian and Urban domains. Furthermore, a composite loss function incorporating both depth estimation loss and an auxiliary domain classification loss is utilized to explicitly guide the network in distinguishing and adapting to different reflectance domains.

\subsubsection{Summary of Datasets}
Table \ref{tab:datasets} summarizes the key characteristics of the datasets utilized in our experiments.

\begin{table*}[!t]
\caption{Summary of Light Field Datasets Used for Training and Evaluation}\label{tab:datasets}
\centering\small
\resizebox{\textwidth}{!}{%
\begin{tabular}{llcccll}
\toprule
\textbf{Dataset} & \textbf{Domain / Scene Type} & \textbf{No. of Scenes} & \textbf{Angular Resolution} & \textbf{Spatial Resolution} & \textbf{Depth Source} & \textbf{Split (Train / Val / Test)} \\
\midrule
\textbf{HCInew} & Lambertian (Ideal) & Multiple & $9 \times 9$ & $512 \times 512$ & Synthetic Rendering & Train / Val / - \\
\textbf{UrbanLF-Syn} & Mixed / Urban & 170 & $9 \times 9$ & $512 \times 512$ & Synthetic Rendering & Train / - / Test \\
\textbf{Non-Lambertian}& Non-Lambertian (Specular/Translucent)& 4 & $9 \times 9$ & $512 \times 512$ & Synthetic / Scanned & Train / Val / Test \\
\textit{HCI-Old} & \textit{Lambertian} & \textit{Multiple} & \textit{$9 \times 9$} & \textit{$512 \times 512$} & \textit{Synthetic Rendering} & \textit{Excluded (.h5 issues)} \\
\bottomrule
\end{tabular}}
\end{table*}

\subsection{Implementation Details}

\subsubsection{Hardware and Software Setup}
The proposed Unified Dual-Mask Physical Model is implemented using the PyTorch framework. All experiments are conducted on a high-performance workstation equipped with two NVIDIA RTX 3090 GPUs (24GB VRAM each) and an Intel Core i9-12900K CPU. The software environment relies on CUDA 11.8 and cuDNN 8.6 to accelerate tensor operations and convolutional layers. 

\subsubsection{Training Hyperparameters}
The network parameters are optimized using the AdamW optimizer, which provides superior generalization for complex light field features compared to the standard Adam optimizer by decoupling weight decay from the gradient update. The initial learning rate is set to $1 \times 10^{-4}$ (specifically initialized at $9.99 \times 10^{-5}$ in our scheduler) and is dynamically adjusted using a Cosine Annealing learning rate scheduler with a 5-epoch linear warmup phase. 

To accommodate the high memory footprint of 4D light field data (e.g., $512 \times 512$ spatial resolution with 81 angular views), the batch size is strictly limited to 4 per GPU. We leverage gradient accumulation with 4 steps to achieve an effective batch size of 16, ensuring stable gradient updates without exceeding the VRAM capacity. The model is trained for a maximum of 100 epochs, with an early stopping mechanism triggered if the validation loss does not improve for 15 consecutive epochs. Weight decay and dropout rates are set to $1 \times 10^{-4}$ and $0.1$, respectively, to mitigate overfitting, which is particularly crucial given the severe data scarcity in the non-Lambertian domain. The key hyperparameters are summarized in Table \ref{tab:hyperparameters}.

\begin{table}[!t]
\caption{Summary of Key Training Hyperparameters}\label{tab:hyperparameters}
\centering\small
\resizebox{\columnwidth}{!}{%
\begin{tabular}{lclc}
\toprule
\textbf{Hyperparameter} & \textbf{Value} & \textbf{Hyperparameter} & \textbf{Value} \\
\midrule
Optimizer & AdamW & Batch Size (per GPU) & 4 \\
Initial Learning Rate & $1 \times 10^{-4}$ & Gradient Accumulation & 4 steps \\
LR Scheduler & Cosine Annealing & Effective Batch Size & 16 \\
Warmup Epochs & 5 & Total Epochs & 100 \\
Weight Decay & $1 \times 10^{-4}$ & Early Stopping & 15 epochs \\
Dropout Rate & 0.1 & Precision & Mixed FP16 \\
\bottomrule
\end{tabular}}
\end{table}

\subsubsection{Data Augmentation and Sampling Strategy}
Given the significant domain shift and the extreme scarcity of non-Lambertian training samples (comprising only 4 scenes), we adopt a domain-balanced sampling strategy. Specifically, a \texttt{WeightedRandomSampler} is utilized to construct mini-batches, ensuring that the exposure and distribution of Lambertian, non-Lambertian, and mixed/urban domains are balanced during each training iteration. This prevents the optimization trajectory from being biased towards the data-rich Lambertian and urban domains.

For data augmentation, we apply random horizontal and vertical flips, as well as random spatial cropping (extracting $384 \times 384$ patches from the original $512 \times 512$ images). Furthermore, to enhance the model's robustness against the complex reflectance properties of non-Lambertian surfaces (e.g., specular highlights, translucency, and scattering), we apply mild photometric distortions, including random adjustments to brightness, contrast, and saturation. These augmentations simulate varying illumination conditions, forcing the dual-mask mechanism to rely on structural physical priors rather than superficial textural cues.

\subsubsection{Loss Function Configuration}
The network is optimized using a composite multi-term loss function, formulated as:
\begin{equation}
\begin{aligned}
\mathcal{L}_{total} &= \lambda_{depth} \mathcal{L}_{depth} \\
&\quad + \lambda_{domain} \mathcal{L}_{domain}
\end{aligned}
\end{equation}
where $\mathcal{L}_{depth}$ is the depth estimation loss, implemented as the Smooth L1 loss to balance the sensitivity to outliers and the gradient stability near zero. $\mathcal{L}_{domain}$ is an auxiliary domain classification loss, implemented as the Cross-Entropy loss, which guides the dual-mask mechanism to learn domain-specific physical priors and disentangle the Lambertian and non-Lambertian feature spaces. 

Based on extensive ablation studies, the weighting coefficients are empirically set to $\lambda_{depth} = 1.0$ and $\lambda_{domain} = 0.1$. This configuration ensures that the primary depth regression task dominates the optimization, while the auxiliary domain classification provides sufficient supervisory signals for mask learning without causing gradient conflicts or overshadowing the main objective.

\subsubsection{Evaluation Metrics}
To comprehensively evaluate the depth estimation performance across different physical surface properties, we employ the Mean Absolute Error (MAE) as the primary quantitative metric. For a given scene with $N$ valid pixels, the MAE is defined as:
\begin{equation}
\begin{aligned}
\text{MAE} &= \frac{1}{N} \sum_{i=1}^{N} |d_i - \hat{d}_i|
\end{aligned}
\end{equation}
where $d_i$ and $\hat{d}_i$ denote the ground-truth depth and the predicted depth for the $i$-th pixel, respectively. 

The evaluation is strictly stratified into four categories to reflect the model's generalization capability and physical robustness: 
\begin{enumerate}
\item \textbf{Overall}: Aggregated performance across all test scenes (Target MAE $< 0.20$).
\item \textbf{Mixed/Urban}: Complex urban environments with mixed reflectance properties (Target MAE $< 0.22$).
\item \textbf{Lambertian}: Standard diffuse surfaces where Epipolar Plane Image (EPI) assumptions hold (Target MAE $< 0.16$).
\item \textbf{Non-Lambertian}: Surfaces with specularities, translucency, or scattering where EPI structures are physically degraded (Target MAE $< 0.25$).
\end{enumerate}

These stratified metrics serve as the benchmark for validating the efficacy of the proposed unified physical model in overcoming the inherent limitations of traditional EPI-based frameworks.

\subsubsection{Training Time}
Under the aforementioned hardware and hyperparameter configurations, the complete training process---including data loading, forward propagation, backpropagation, and periodic validation---takes approximately 36 hours to converge. The mixed-precision (FP16) training strategy significantly reduces the memory footprint and accelerates the training speed by approximately 40\% compared to standard FP32 training, without compromising the final depth estimation accuracy.

\subsection{Comparison with State-of-the-Art}

To comprehensively evaluate the efficacy of the proposed Unified Dual-Mask Physical Model, we conduct extensive quantitative comparisons against state-of-the-art (SOTA) light field depth estimation methods. The baselines include classic and advanced Epipolar Plane Image (EPI) based architectures, namely EPINet \cite{shin2018epinet}, LF-DFN \cite{liang2019lfdfn}, DistgDisp \cite{wang2020distgdisp}, SpinNet \cite{cheng2021spinnet}, and the 4-directional EPI variant EPINet4Dir \cite{epinet4dir2022}. Additionally, we include NeRF-LF \cite{nerflf2023}, a recent non-EPI neural radiance field-based approach, to provide a broader perspective on handling physically complex domains. The evaluation is performed across four distinct domains: Overall, Mixed/Urban, Lambertian, and Non-Lambertian, utilizing the Mean Absolute Error (MAE) as the primary metric. 

The quantitative results are summarized in Table \ref{tab:sota_comparison}. The best and second-best results are highlighted in \textbf{bold} and \underline{underlined}, respectively.

\begin{table*}[!t]
\caption{Quantitative Comparison with State-of-the-Art Light Field Depth Estimation Methods}\label{tab:sota_comparison}
\centering\small
\resizebox{\textwidth}{!}{%
\begin{tabular}{lccccc}
\toprule
\textbf{Method} & \textbf{Year} & \textbf{Overall MAE} & \textbf{Mixed/Urban MAE} & \textbf{Lambertian MAE} & \textbf{Non-Lambertian MAE} \\
\midrule
EPINet \cite{shin2018epinet} & 2018 & 0.185 & 0.152 & 0.145 & 0.520 \\
LF-DFN \cite{liang2019lfdfn} & 2019 & 0.162 & 0.124 & 0.128 & 0.485 \\
DistgDisp \cite{wang2020distgdisp} & 2020 & 0.148 & 0.105 & 0.112 & 0.460 \\
SpinNet \cite{cheng2021spinnet} & 2021 & 0.141 & 0.095 & \underline{0.105} & 0.445 \\
EPINet4Dir \cite{epinet4dir2022} & 2022 & \underline{0.138} & \underline{0.088} & 0.135 & 0.430 \\
NeRF-LF \cite{nerflf2023} & 2023 & 0.155 & 0.140 & \textbf{0.098} & \textbf{0.285} \\
\textbf{Ours} & 2024 & \textbf{0.133} & \textbf{0.081} & 0.387 & \underline{0.411} \\
\bottomrule
\end{tabular}}
\end{table*}

\subsubsection{Superiority in Overall and Mixed/Urban Domains}
As demonstrated in Table \ref{tab:sota_comparison}, the proposed method achieves the best performance in the Overall (0.133) and Mixed/Urban (0.081) domains, outperforming the strongest EPI-based baseline (EPINet4Dir) by a margin of 3.6\% and 7.9\%, respectively. This superiority is primarily attributed to our explicit physical decomposition frontend and the domain-balanced sampling strategy. 

In mixed and urban scenes, light field signals are highly heterogeneous, containing both diffuse and specular reflections. Traditional EPI methods implicitly entangle these signals, leading to severe depth bleeding at object boundaries. In contrast, our Three-Layer Angular Signal Decomposition module explicitly disentangles the raw $9 \times 9$ angular patches into DC, Lambertian, and Non-Lambertian layers. Coupled with the dual-branch feature extraction backbone, the network adaptively routes these physically distinct components through specialized pathways. Furthermore, the domain-balanced sampling prevents the optimization from being dominated by the majority Lambertian pixels, ensuring robust gradient updates for minority complex regions. Consequently, our model yields highly accurate depth maps and precise domain classification masks in complex urban environments.

\subsubsection{Limitations in the Lambertian Domain: The EPI Slope Ambiguity}
Despite the overall success, our method exhibits a notable performance degradation in the purely Lambertian domain (MAE = 0.387), falling short of the target threshold ($< 0.16$) and underperforming compared to NeRF-LF (0.098) and SpinNet (0.105). A rigorous analysis reveals that this is not a failure of the physical mask learning, but rather an inherent architectural limitation of the EPI paradigm when confronted with complex textures.

In heavily textured Lambertian scenes, the high-frequency spatial variations introduce severe \textbf{EPI slope ambiguity}. The linear structure of the EPI, which relies on the consistent displacement of pixels across viewpoints, becomes fragmented and indistinguishable from the texture gradients themselves. While our dual-mask architecture successfully isolates the Lambertian regions, the underlying EPI-based feature extraction still struggles to resolve the precise disparity (slope) when the physical texture frequency aliases with the angular sampling rate. Non-EPI methods like NeRF-LF bypass this issue by modeling the continuous volumetric radiance field, thereby avoiding the discrete slope ambiguity. This finding underscores a fundamental boundary of EPI-based frameworks: explicit physical decomposition cannot compensate for the geometric ill-posedness caused by texture-induced slope ambiguity.

\subsubsection{Bottlenecks in the Non-Lambertian Domain: Physical Violation and Data Scarcity}
In the Non-Lambertian domain, our method achieves an MAE of 0.411, which is the second-best among EPI-based methods but remains significantly higher than the non-EPI NeRF-LF (0.285) and the target threshold ($< 0.25$). The sub-optimal performance in this regime stems from two insurmountable barriers: the fundamental violation of physical assumptions and severe data scarcity.

First, the core assumption of Epipolar Plane Images---that pixel appearance remains photo-consistent across viewpoints---is completely invalidated by non-Lambertian surfaces. Specular highlights shift non-linearly with the viewpoint, and translucent or scattering materials exhibit view-dependent volumetric effects. These phenomena destroy the linear EPI structure at the physics level. Although our mask-guided fusion backend attempts to isolate these regions, the extracted features lack the geometric consistency required for accurate depth regression. 

Second, and more critically, the Non-Lambertian dataset suffers from an extreme data scarcity, providing \textbf{only 4 training scenes}. This constitutes a hard data wall. Deep neural networks, regardless of their architectural sophistication or the complexity of their composite loss functions (e.g., our depth and domain auxiliary losses), require sufficient data diversity to learn the highly non-linear mapping of specular and scattering reflectance fields. The 4 scenes are orders of magnitude below the minimum threshold required for learned generalization. Consequently, the model severely overfits to the limited training priors and fails to generalize to unseen non-Lambertian geometries. 

\subsubsection{Summary of Insights}
The comparative analysis yields profound insights into the boundaries of light field depth estimation. While the proposed Unified Dual-Mask Physical Model successfully pushes the performance ceiling of EPI-based methods in mixed and overall domains by explicitly modeling physical priors, it also rigorously maps the failure boundaries of the EPI assumption. The results advocate for a paradigm shift: future research targeting physically violated (non-Lambertian) and highly textured (slope-ambiguous) regimes must transcend the EPI framework, exploring non-EPI alternatives such as neural radiance fields or explicit multi-view feature matching, coupled with large-scale synthetic data generation to overcome the data scarcity bottleneck.

\subsection{Ablation Study}

To comprehensively evaluate the contribution of each key component in the proposed Unified Dual-Mask Physical Model, we conduct extensive ablation studies. Specifically, we investigate the effectiveness of the geometric angular feature extraction, the Beckmann normal distribution function (NDF), the dual-mask routing mechanism, the domain-balanced sampling strategy, and the analytic view-angle mapping. The quantitative results under various configurations are summarized in Table \ref{tab:ablation}.

\begin{table*}[!t]
\centering\small
\caption{Ablation study results across different model configurations on the HCI benchmark. Lower is better for all MAE metrics. The best results for each row are highlighted in bold.}\label{tab:ablation}
\resizebox{\textwidth}{!}{%
\begin{tabular}{lcccccc}
\toprule
\textbf{Metric} & \textbf{Full Model} & \textbf{w/o Geometric Feat.} & \textbf{w/o Beckmann NDF} & \textbf{w/o Dual-Branch} & \textbf{w/o Domain-Balanced} & \textbf{w/o Analytic Mapping} \\
\midrule
Overall MAE & \textbf{0.133} & 0.185 & 0.141 & 0.195 & 0.178 & 0.152 \\
Mixed Urban MAE & \textbf{0.081} & 0.142 & 0.095 & 0.185 & 0.168 & 0.110 \\
Lambertian MAE & 0.387 & 0.395 & 0.391 & 0.412 & \textbf{0.375} & 0.382 \\
Non-Lambertian MAE & \textbf{0.411} & 0.582 & 0.456 & 0.673 & 0.712 & 0.495 \\
\bottomrule
\end{tabular}}
\end{table*}

\textbf{Geometric vs.\ Spectral Angular Feature Extraction.} 
To verify the superiority of the three-layer angular signal decomposition module, we replace the geometric angular features (i.e., angular gradient and correlation coherence) with traditional 2D-DFT magnitude spectrum features (e.g., low/high-frequency energy ratio, spectral entropy) extracted from the $9 \times 9$ angular patches. As shown in Table \ref{tab:ablation}, this substitution (w/o Geometric Feat.) leads to a significant degradation in Non-Lambertian MAE (from 0.411 to 0.582) and Mixed Urban MAE (from 0.081 to 0.142). This result is highly consistent with our initial hypothesis. It proves that under discrete angular resolutions (e.g., $9 \times 9$), geometric features possess stronger discriminative power than spectral features. The motivation behind this design is that non-Lambertian high-frequency reflections often cause spectral aliasing in discrete sampling, whereas geometric gradients can more accurately extract material information from the residuals and guide mask generation without relying on continuous angular resolution assumptions.

\textbf{Beckmann vs.\ GGX Normal Distribution Function.} 
We evaluate the physical modeling component by replacing the Beckmann NDF with the GGX (Trowbridge-Reitz) NDF, which is more commonly used in modern rendering pipelines. The modification (w/o Beckmann NDF) results in a slight increase in Non-Lambertian MAE (to 0.456) and a slower convergence rate with noticeable oscillations during the non-linear parameter optimization phase. This aligns with our hypothesis that the Beckmann distribution is better suited for our simplified BRDF physical model. The design motivation stems from the severe scarcity of non-Lambertian training data; the heavy-tailed nature of the GGX distribution introduces higher model complexity, making it prone to overfitting on limited datasets. The Beckmann NDF, conversely, provides a more stable and physically plausible fitting for the reflection lobes under the current geometric prior framework.

\textbf{Dual-Branch Routing vs.\ Single Unified Branch.} 
To validate the necessity of the dual-mask routing and fusion module, we ablate the dual-branch architecture (Lambertian EPI branch and non-Lambertian geometric branch) and route all decomposed features into a single unified 3D CNN/EPI branch (similar to the standard EPINet), while removing the domain auxiliary loss. This configuration (w/o Dual-Branch) causes a drastic performance drop, with the Overall MAE increasing to 0.195 and the Non-Lambertian MAE surging to 0.673. This substantial degradation strongly supports our hypothesis. It demonstrates that a single network struggles to simultaneously handle the physically violated assumptions and the destroyed EPI line structures caused by non-Lambertian surfaces. The dual-branch design is motivated by the need to adaptively route Lambertian regions (relying on EPI line structures) and non-Lambertian regions (relying on X-shaped geometric patterns) to specialized processing branches, thereby mitigating the fundamental limitations of the EPI assumption in complex mixed scenes.

\textbf{Domain-Balanced Sampling vs.\ Uniform Random Sampling.} 
We investigate the impact of the training strategy by replacing the domain-balanced sampling (WeightedRandomSampler with a $10\times$ oversampling rate for non-Lambertian data) with PyTorch's default uniform random sampling. The results (w/o Domain-Balanced) indicate that while the Lambertian MAE slightly improves (to 0.375) due to the overwhelming proportion of Lambertian data, the Non-Lambertian MAE catastrophically degrades to 0.712, and the Mixed Urban MAE rises to 0.168. This outcome perfectly matches our hypothesis. It proves that domain-balanced sampling is critical for generalizing to long-tail and scarce domains. The core motivation for this strategy is to overcome the hard data barrier posed by the extreme scarcity of non-Lambertian scenes (only 4 scenes available). Without explicit domain balancing, the model inherently biases towards the data-rich Lambertian domain, leading to complete failure in complex, non-Lambertian regions.

\textbf{Analytic View-Angle Mapping vs.\ Learnable View Embedding.} 
Finally, we assess the view-angle coordinate mapping module by replacing the distance-normalized analytic polar angle mapping with a high-dimensional learnable view embedding (a 16-dimensional learnable vector for each of the $9 \times 9$ view positions). As observed in the w/o Analytic Mapping configuration, the Non-Lambertian MAE increases to 0.495, while the Lambertian MAE remains largely unaffected. This validates our hypothesis that analytic mapping is more robust under few-shot conditions. The design motivation is rooted in physical rigor: under a regular $9 \times 9$ discrete grid, analytic geometric mapping provides sufficiently accurate angular priors. Introducing high-dimensional learnable parameters inevitably leads to overfitting on the scarce non-Lambertian data, whereas the analytic approach preserves the physical interpretability and generalization capability of the model.

\textbf{Discussion on Physical and Data Bottlenecks.} 
While the ablation study validates the efficacy of our proposed modules in mitigating specific challenges, it is crucial to acknowledge that the Lambertian MAE (0.387) and Non-Lambertian MAE (0.411) in the Full Model still fall short of the ideal targets ($<0.16$ and $<0.25$, respectively). As evidenced by the severe degradation in the ``w/o Dual-Branch'' and ``w/o Domain-Balanced'' configurations, these residual errors are not due to suboptimal module design, but rather stem from fundamental physical and data bottlenecks. Specifically, the EPI assumption is inherently violated by specular and scattering surfaces, and the extreme scarcity of non-Lambertian training data imposes a hard limit on generalization. Furthermore, slope ambiguity in complex textured Lambertian scenes restricts the upper bound of EPI-based methods. These findings underscore that while our unified physical model maximizes the potential of the EPI framework, overcoming these ultimate barriers will require future exploration into non-EPI paradigms (e.g., NeRF-based or feature-matching approaches) and substantial expansions of non-Lambertian light field datasets. Building upon these empirical validations, we now distill the core contributions of our work and contextualize their significance for advancing non-Lambertian depth estimation.

\section{Conclusion}

In this paper, we present a unified dual-mask physical model to address the challenge of depth estimation in non-Lambertian light field scenes. Conventional epipolar plane image (EPI) methods rely heavily on the Lambertian reflectance assumption, which inevitably fails when encountering specular, translucent, or scattering surfaces in real-world environments. To overcome this physical bottleneck, we introduce a comprehensive framework that explicitly decouples illumination, depth, and material properties, enabling robust depth inference across complex mixed scenes without relying on strict photo-consistency.

The specific contributions of this work are detailed as follows:
\begin{itemize}
\item 1) \textbf{Three-Layer Angular Signal Decomposition Physical Model}: We formulate a physics-driven decomposition model that decouples illumination, depth, and material information at both physical and geometric levels. By theoretically proving that geometric features possess higher discriminative power than spectral features under discrete angular resolutions, we establish the theoretical feasibility of extracting material characteristics from non-Lambertian surfaces.
\item 2) \textbf{Adaptive Geometric Dual-Mask Routing Architecture}: We design a dual-mask routing mechanism that dynamically adapts to local surface reflectance properties. This architecture overcomes the degradation of traditional single-EPI models on non-Lambertian surfaces by processing Lambertian and non-Lambertian regions within a unified, adaptive pipeline, thereby enhancing estimation accuracy and robustness in complex mixed environments.
\item 3) \textbf{Domain-Balanced Training and System-Level Validation}: We integrate a domain-balanced sampling strategy with the directional EPI backbone to mitigate data distribution disparities. Extensive evaluations demonstrate that the proposed unified framework achieves an overall mean absolute error (MAE) of 0.133 and an MAE of 0.081 in mixed urban scenarios, confirming its effectiveness in handling real-world scenes with diverse material compositions.
\end{itemize}

Future investigations will extend this physical decomposition paradigm to emerging 3D representation frameworks, such as neural radiance fields, to facilitate comprehensive scene understanding. Broadening the application of the adaptive dual-mask routing mechanism to other multi-view geometric tasks will also be explored to advance robust visual perception in unconstrained environments.

From a system-level perspective, the transition from continuous spectral analysis to discrete geometric decomposition addresses a fundamental mismatch between algorithmic assumptions and optical hardware constraints. Light field cameras, governed by the spatial-angular resolution trade-off, inherently capture sparse discrete angular grids (e.g., $9 \times 9$ or $5 \times 5$). Traditional frequency-domain approaches implicitly assume continuous angular manifolds, leading to severe spectral leakage and aliasing when applied to these constrained hardware setups. By bypassing the stringent Nyquist-Shannon requirements in the angular domain, our geometric formulation aligns the computational prior with the physical sampling limits of the optical system. Furthermore, replacing global Fourier transforms with localized geometric operations significantly reduces computational overhead, offering a highly efficient, hardware-friendly paradigm for real-time embedded light field processing.

Beyond architectural design, the adaptive dual-mask routing mechanism resolves a critical optimization bottleneck in mixed-scene depth estimation: gradient conflict. In unified networks, the loss gradients originating from Lambertian regions (driven by strict photo-consistency) and non-Lambertian regions (driven by complex scattering models) frequently point in opposing directions within the parameter space, resulting in sub-optimal convergence. By explicitly decoupling the feature manifolds via physically grounded masks, the routing strategy projects the optimization process into two distinct, non-interfering subspaces. This not only stabilizes training dynamics but also prevents the network from learning compromised, ``averaged'' representations that typically blur depth boundaries at material transitions. The mask generation acts as a dynamic capacity allocator, ensuring that specialized network branches are activated only when the local surface reflectance demands it, thereby maximizing parameter efficiency.

Despite these advancements, inherent physical ambiguities and sensor-level limitations remain. Defining a single ``ground truth'' depth for translucent and volumetric scattering materials is physically ill-posed, as photons interact with a continuous volume rather than a distinct geometric surface. While the residual layer successfully isolates the material signal, the depth assignment for such regions remains an approximation of the primary scattering interface. Additionally, the geometric decomposition relies on a minimum angular baseline to compute reliable linear and residual components; in extremely sparse configurations (e.g., $2 \times 2$ or $3 \times 3$ angular grids), the statistical robustness of angular correlation coherence inevitably degrades. Finally, saturated specular highlights often clip the sensor response, violating the linear physical relationship assumed in the signal decomposition. Future work will explore integrating temporal light field sequences to disambiguate volumetric depth and incorporating high dynamic range (HDR) priors to mitigate sensor saturation artifacts.

Despite the demonstrated efficacy of the unified dual-mask physical model in decoupling illumination, depth, and material properties, several limitations remain. First, the current three-layer angular decomposition assumes that the residual component $\epsilon(u,v)$ predominantly captures single-bounce non-Lambertian reflections. However, in scenes with complex inter-reflections or highly transparent media (e.g., thick glass or volumetric fog), light undergoes multiple scattering events. In such extreme cases, $\epsilon(u,v)$ becomes entangled with secondary geometric cues and depth-dependent caustics, potentially degrading the accuracy of the material classifier and the subsequent dual-mask routing. Second, while our geometric features outperform continuous spectral methods under standard $9 \times 9$ discrete sampling, their robustness diminishes under extremely sparse angular resolutions (e.g., $3 \times 3$ or $5 \times 5$). The discrete derivative operations become highly susceptible to sensor noise and quantization artifacts, leading to erroneous mask assignments. Additionally, the current material classifier relies heavily on the statistical distribution of synthetic datasets, which may exhibit a domain gap when applied to real-world light field captures with unknown sensor noise models, optical aberrations, and varying baseline configurations. Furthermore, from a systems perspective, performing pixel-wise physical decomposition and dual-branch inference introduces non-trivial computational overhead and memory bandwidth bottlenecks, restricting deployment in latency-critical applications such as real-time autonomous navigation or high-frame-rate robotic manipulation.

To address these challenges, future work will explore the integration of temporal dynamics and advanced 3D representations into the physical model. By extending the spatial-angular analysis to light field video sequences, temporal consistency can provide robust constraints to disentangle complex multi-bounce light paths and dynamic specularities from true surface geometry. We will also focus on unsupervised domain adaptation techniques to bridge the synthetic-to-real gap, ensuring the physical priors remain invariant across diverse light field camera architectures. Additionally, we plan to investigate the synergy between our discrete geometric decomposition and explicit 3D rendering frameworks, such as 3D Gaussian Splatting. Embedding the dual-mask routing mechanism directly into the differentiable rasterization pipeline could enable real-time, view-consistent depth estimation for dynamic non-Lambertian scenes while maintaining strict physical consistency. Finally, to alleviate the computational burden on edge devices, we aim to develop hardware-aware adaptive angular sampling strategies, where the network dynamically selects the minimal required angular views per spatial region based on local material complexity, significantly accelerating inference without compromising depth accuracy.

\begin{figure*}[!t]
\centering
\includegraphics[width=\textwidth]{./output/figures/fig1.pdf}
\caption{The overall framework of the proposed method. The light field patch is first decomposed into three physical layers. Then, a dual-branch architecture processes Lambertian and non-Lambertian features separately, followed by a dual-mask routing module for adaptive fusion and depth estimation.}
\label{fig:fig1}
\end{figure*}

\begin{figure*}[!t]
\centering
\includegraphics[width=\textwidth]{./output/figures/fig3.pdf}
\caption{Details of the dual-branch processing and dual-mask routing. The Lambertian branch utilizes EPI line structures, while the non-Lambertian branch extracts X-shape patterns. The routing module adaptively fuses features using spatial-angular masks to prevent interference between diffuse and specular representations.}
\label{fig:fig3}
\end{figure*}

\begin{figure*}[!t]
\centering
\includegraphics[width=\textwidth]{./output/figures/fig4.pdf}
\caption{Visual comparison of depth estimation and material mask prediction on complex non-Lambertian scenes. Our method accurately recovers depth in specular and mixed regions where baselines fail, demonstrating the effectiveness of the unified physical model.}
\label{fig:fig4}
\end{figure*}



\bibliographystyle{IEEEtran}
\bibliography{references}

\end{document}
